\sectioncentered*{Заключение}
\addcontentsline{toc}{section}{Заключение}

В результате выполнения дипломного проекта получены следующие результаты:

\begin{enumerate}[label=\arabic*]
    \item В аналитическом разделе исследованы особенности организации времени обучающихся, рассмотрены факторы академической прокрастинации и выполнен обзор существующих цифровых решений в области планирования, управления обучением и интеллектуального ассистирования. Проведенный анализ показал, что универсальные инструменты не обеспечивают достаточной поддержки декомпозиции долгосрочных учебных задач, учета академического контекста и последующего переноса результата в персональный календарный план.
    \item В разделе проектирования сформулированы функциональные и нефункциональные требования к персональному ассистенту организации времени обучающихся. Разработаны функциональная модель, логическая архитектура, модель данных и ключевые сценарии взаимодействия пользователя с системой. Обосновано применение архитектуры модульного монолита с разделением на слои контроллеров, сервисов и репозиториев, а также использование агентного модуля, в котором языковая модель выполняет смысловую декомпозицию пользовательского запроса, а детерминированный алгоритм осуществляет распределение подзадач по временным интервалам календаря.
    \item В разделе реализации разработан программный комплекс персонального ассистента организации времени обучающихся. Реализованы серверная часть на базе языка Java и фреймворка Spring Boot, подсистема хранения данных на основе СУБД PostgreSQL, агентный модуль с поддержкой структурированного вывода языковой модели, механизм формирования, подтверждения и отклонения черновика расписания, а также клиентское одностраничное приложение с календарной визуализацией результата. Выполнено функциональное тестирование, которое продемонстрировало полноту и корректность реализованных функций.
    \item В экономическом разделе выполнено технико-экономическое обоснование разработки и использования программного средства. Расчеты показали экономическую целесообразность проекта как для организации-разработчика, так и для организации-заказчика.
\end{enumerate}

Практическая значимость результатов состоит в создании программного инструмента, который может быть использован для повышения прозрачности работы над долгосрочными учебными задачами, снижения когнитивной нагрузки и поддержки навыков самоорганизации обучающихся. Разработанная система предполагается к использованию в качестве программного средства поддержки учебного планирования и организации времени обучающихся для персональных нужд и внедрения в цифровую среду высших учебных заведений.

Дальнейшее развитие работы может быть связано с расширением источников учебного контекста, совершенствованием механизмов персонализации, развитием интерфейсных средств сопровождения пользователя и углублением интеграции с цифровой образовательной средой.
